\exercisesection{6}

¶ 1. p-进域  的每个元素 z 都可唯一写成 \(Q_{p}\)，其中 \(p^{h}\sum_{k=0}^{\infty}c_{k}p^{k}\)、\(h\in Z\)，且对所有 \(c_{0}\neq0\) 有 （z 的“p-进展开”）。\(0\leqslant c_{k}<p\)\(k\geqslant0\)（z 的“p-进展开”）。

\begin{enumerate}
\def\labelenumi{(\alph{enumi})}
\tightlist
\item
  证明  的充要条件是存在两个整数 \(z \in \mathbf{Q}\)，使得对所有 \(m \geqslant 0, n \geqslant 1\) 都有 。（注意，若 \(c_{k+n} = c_k\)，其中 b 不被 p 整除，则\(k \geqslant m\)\(z = a / b \in \mathbf{Q}\)\(b\)\(p\)，其中 b 不被 p 整除，则
\end{enumerate}

\[
a - b \sum_ {k = 0} ^ {n} c _ {k} p ^ {k} = a _ {n + 1} p ^ {n + 1}
\]

其中 \(a_{n+1}\) 为整数，且序列 \(|a_k|\)（通常绝对值）有界。）

\begin{enumerate}
\def\labelenumi{(\alph{enumi})}
\setcounter{enumi}{1}
\tightlist
\item
  假设满足  的整数 k 的递增序列  满足 \((k_n)\)。证明 z 在 \(k\)\(c_k \neq 0\)\(\limsup_{n \to \infty} (k_{n+1}/k_n) = +\infty\) 上超越。（对 \(z\)\(\mathbf{Q}\)，令 \(x \in \mathbf{Q}\) 与 \(|x|\) 分别表示通常绝对值与 p-进绝对值。假设 \(|x|_p\) 是满足 \(p\)\(\mathrm{P} \in \mathbf{Z}[\mathrm{X}]\) 的不可约多项式；若 \(\mathrm{P}(z) = 0\)，使用 Taylor 公式证明 \(z_n = p^n \sum_{k=0}^{\infty} c_k p^k\) 随 n 趋于 \(|\mathrm{P}(z_n)/(z - z_n)|_p\) 时趋于一个  的极限；然后考虑通常绝对值 \(\neq 0\)，从假设得到矛盾。）\(n\)\(+\infty\)\(|\mathrm{P}(z_n)|\)，从假设得到矛盾。）
\end{enumerate}

\begin{enumerate}
\def\labelenumi{\arabic{enumi}.}
\setcounter{enumi}{1}
\tightlist
\item
  设 D 为特征 \(\neq2\) 的非交换域，Z 为其中心。假设 D 中每个包含 Z 的交换子域 K 在 Z 上的次数均 \(\leqslant2\)。
\end{enumerate}

\begin{enumerate}
\def\labelenumi{(\alph{enumi})}
\item
  不使用引理 1，证明 。（取 \([\mathbf{D}:\mathbf{Z}] = 4\) 使得 \(a\in \mathbf{D} - \mathbf{Z}\)，并对所有 \(a^2\in \mathbf{Z}\) 令 ；注意 \(\sigma (x) = axa^{-1}\) 分解为 Z 上两个向量子空间 \(x\in \mathbf{D}\)\(\mathbf{D}\) 与 \(\mathbf{Z}\)\(\mathrm{D}_{+}\) 的直和，使得 \(\mathrm{D}_{-}\) 在 \(\sigma (x) = x\) 上成立，而 \(\mathrm{D}_{+}\) 在 \(\sigma (x) = -x\) 上成立；还注意 \(\mathrm{D}_{-}\) 是 \(\mathrm{D}_{+}\) 的子域，而 \(\mathbf{D}\) 是 \(\mathrm{D}_{-}\) 上的 1 维向量空间；最后证明 \(\mathrm{D}_{+}\) 不可能不同于 \(\mathrm{D}_{+}\)。）\(\mathbf{Z}(a)\)。）
\item
  证明  是 Z 上的四元数代数。（借助 (a) 构造 \(\mathbf{D}\) 在 Z 上的一组基。）\(\mathbf{Z}\)\(\mathbf{D}\)\(\mathbf{Z}\) 在 Z 上的一组基。）
\end{enumerate}

